% ============================================================
% Section 123: 半金属交叠能带与费米能级
% ============================================================
\section{固体物理：半金属交叠能带与费米能级}
\label{sec:semimetal-overlap}

\begin{modelbox}[题目：半金属能带交叠时的费米能级位置]
半金属交叠的能带为
\[
E_1(k)=E_1(0)-\frac{\hbar^2k^2}{2m_1},\qquad m_1=0.18m,
\]
\[
E_2(k)=E_2(k_0)+\frac{\hbar^2}{2m_2}(k-k_0)^2,\qquad m_2=0.06m,
\]
其中 $E_1(0)$ 为带 1 的带顶，$E_2(k_0)$ 为带 2 的带底。交叠部分 $E_1(0)-E_2(k_0)=0.1\,\text{eV}$。由于能带交叠，能带 1 中的部分电子转移到带 2，而在带 1 中形成空穴。讨论 $T=0\,\text{K}$ 时的费米能级位置。
\end{modelbox}

\begin{tipbox}[考点快速分析]
\begin{itemize}
    \item \textbf{半金属特征}：导带和价带发生微小交叠，$T=0$ 时即存在电子和空穴
    \item \textbf{电中性条件}：$n=p$，转移的电子数等于产生的空穴数
    \item \textbf{态密度}：三维自由电子型 $g(E)\propto m^{3/2}E^{1/2}$
    \item \textbf{费米能级平衡}：由两带有效质量（态密度权重）共同决定
\end{itemize}
\end{tipbox}

\begin{derivationbox}[标准解题过程]
\textbf{1. 物理模型与能带交叠}

带 1 为空穴带（价带），开口向下，带顶在 $E_1(0)$，空穴有效质量 $m_1=0.18m$。

带 2 为电子带（导带），开口向上，带底在 $E_2(k_0)$，电子有效质量 $m_2=0.06m$。

两带发生交叠，交叠量 $\Delta=E_1(0)-E_2(k_0)=0.1\,\text{eV}$。由于带 1 顶部能量高于带 2 底部，带 1 顶部的电子自发转移到带 2 底部，在带 1 中留下等量的空穴。系统达到平衡时，费米能级 $E_F$ 位于交叠区内。

设带 2 中费米能级距带底的高度为 $\delta_n$，带 1 中费米能级距带顶的深度为 $\delta_p$，则
\[
\delta_n+\delta_p=\Delta=0.1\,\text{eV}.
\]

\textbf{2. 绝对零度下的载流子浓度}

在 $T=0\,\text{K}$，费米-狄拉克分布退化为阶跃函数。

\textit{带 2 电子浓度}（以带底为零点）：
\[
n=\int_0^{\delta_n}g_2(E)\,dE=\frac{1}{3\pi^2}\left(\frac{2m_2}{\hbar^2}\right)^{3/2}\delta_n^{3/2}\propto m_2^{3/2}\delta_n^{3/2}.
\]

\textit{带 1 空穴浓度}（以带顶为零点，能量向下为正）：
\[
p=\int_0^{\delta_p}g_1(E)\,dE=\frac{1}{3\pi^2}\left(\frac{2m_1}{\hbar^2}\right)^{3/2}\delta_p^{3/2}\propto m_1^{3/2}\delta_p^{3/2}.
\]

\textbf{3. 电中性条件与求解}

系统电中性要求 $n=p$：
\[
m_2^{3/2}\delta_n^{3/2}=m_1^{3/2}\delta_p^{3/2}
\implies m_2\delta_n=m_1\delta_p.
\]

代入 $m_1=0.18m$，$m_2=0.06m$：
\[
0.06\,\delta_n=0.18\,\delta_p
\implies \delta_n=3\delta_p.
\]

联立 $\delta_n+\delta_p=0.1\,\text{eV}$：
\[
3\delta_p+\delta_p=0.1
\implies 4\delta_p=0.1
\implies \delta_p=0.025\,\text{eV},
\]
\[
\delta_n=0.1-0.025=0.075\,\text{eV}.
\]

\textit{结论}：
\begin{itemize}
    \item 费米能级 $E_F$ 位于带 2 底以上 $\delta_n=0.075\,\text{eV}$ 处；
    \item 费米能级位于带 1 顶以下 $\delta_p=0.025\,\text{eV}$ 处。
\end{itemize}

\begin{equation}
\boxed{\delta_n=0.075\,\text{eV},\quad \delta_p=0.025\,\text{eV}.}
\end{equation}

\textit{物理分析}：$m_1>m_2$，带 1（空穴带）的有效质量大，态密度大。为了与带 2 提供等量的载流子，费米能级只需在带 1 中切入较浅的深度（$\delta_p$ 小），而在带 2 中需要填充较高的能量（$\delta_n$ 大）。因此费米能级更靠近有效质量大的带 1（价带顶）。
\end{derivationbox}

\begin{formulabox}[答案汇总]
\begin{center}
\begin{tabular}{ll}
\toprule
\textbf{物理量} & \textbf{数值} \\
\midrule
带 2 电子费米深度 $\delta_n$ & $0.075\,\text{eV}$（带底以上） \\
带 1 空穴费米深度 $\delta_p$ & $0.025\,\text{eV}$（带顶以下） \\
费米能级位置 & 交叠区内，距带 2 底 $0.075\,\text{eV}$，距带 1 顶 $0.025\,\text{eV}$ \\
\bottomrule
\end{tabular}
\end{center}
\end{formulabox}

\begin{insightbox}[物理图像]
\begin{itemize}
    \item 半金属（如铋、锑、石墨）的导带和价带发生微小交叠，导致即使在绝对零度也存在有限数量的电子和空穴。这与本征半导体不同，后者的载流子由热激发产生，$T=0$ 时 $n=p=0$。
    \item 绝对零度下的费米能级由两带态密度权重的平衡决定：有效质量越大（态密度越大），费米能级距该带极值越近。这是因为大有效质量的能带``更容易''容纳载流子。
    \item 半金属的电子和空穴浓度通常很小（$\sim10^{17\text{--}20\,\text{cm}^{-3}$），远低于普通金属（$\sim10^{22}\,\text{cm}^{-3}$），因此电阻率较高，具有独特的输运性质（如巨大的磁阻效应）。
\end{itemize}
\end{insightbox}
